\documentclass[utf8, bachelor]{gradu3}
%DIF LATEXDIFF DIFFERENCE FILE


% Jos työ on kandidaatintutkielma eikä pro gradu, käytä ylläolevan asemesta
%\documentclass[utf8,bachelor]{gradu3}
% Jos kirjoitat englanniksi, käytä ylläolevan asemesta
%\documentclass[utf8,english]{gradu3}
% tai
%\documentclass[utf8,bachelor,english]{gradu3}

\usepackage{graphicx} % kuvien mukaan ottamista varten

\usepackage{amsmath} % hyödyllinen jos tekstisi sisältää matikkaa, 

\usepackage{booktabs} % hyvä kauniiden taulukoiden tekemiseen

% HUOM! Tämän tulee olla viimeinen \usepackage koko dokumentissa!
\usepackage[bookmarksopen,bookmarksnumbered,linktocpage]{hyperref}

\addbibresource{Annan-kandi.bib} % Lähdetietokannan tiedostonimi
%DIF PREAMBLE EXTENSION ADDED BY LATEXDIFF
%DIF UNDERLINE PREAMBLE %DIF PREAMBLE
\RequirePackage[normalem]{ulem} %DIF PREAMBLE
\RequirePackage{color}\definecolor{RED}{rgb}{1,0,0}\definecolor{BLUE}{rgb}{0,0,1} %DIF PREAMBLE
\providecommand{\DIFaddtex}[1]{{\protect\color{blue}\uwave{#1}}} %DIF PREAMBLE
\providecommand{\DIFdeltex}[1]{{\protect\color{red}\sout{#1}}}                      %DIF PREAMBLE
%DIF SAFE PREAMBLE %DIF PREAMBLE
\providecommand{\DIFaddbegin}{} %DIF PREAMBLE
\providecommand{\DIFaddend}{} %DIF PREAMBLE
\providecommand{\DIFdelbegin}{} %DIF PREAMBLE
\providecommand{\DIFdelend}{} %DIF PREAMBLE
\providecommand{\DIFmodbegin}{} %DIF PREAMBLE
\providecommand{\DIFmodend}{} %DIF PREAMBLE
%DIF FLOATSAFE PREAMBLE %DIF PREAMBLE
\providecommand{\DIFaddFL}[1]{\DIFadd{#1}} %DIF PREAMBLE
\providecommand{\DIFdelFL}[1]{\DIFdel{#1}} %DIF PREAMBLE
\providecommand{\DIFaddbeginFL}{} %DIF PREAMBLE
\providecommand{\DIFaddendFL}{} %DIF PREAMBLE
\providecommand{\DIFdelbeginFL}{} %DIF PREAMBLE
\providecommand{\DIFdelendFL}{} %DIF PREAMBLE
%DIF HYPERREF PREAMBLE %DIF PREAMBLE
\providecommand{\DIFadd}[1]{\texorpdfstring{\DIFaddtex{#1}}{#1}} %DIF PREAMBLE
\providecommand{\DIFdel}[1]{\texorpdfstring{\DIFdeltex{#1}}{}} %DIF PREAMBLE
\newcommand{\DIFscaledelfig}{0.5}
%DIF HIGHLIGHTGRAPHICS PREAMBLE %DIF PREAMBLE
\RequirePackage{settobox} %DIF PREAMBLE
\RequirePackage{letltxmacro} %DIF PREAMBLE
\newsavebox{\DIFdelgraphicsbox} %DIF PREAMBLE
\newlength{\DIFdelgraphicswidth} %DIF PREAMBLE
\newlength{\DIFdelgraphicsheight} %DIF PREAMBLE
% store original definition of \includegraphics %DIF PREAMBLE
\LetLtxMacro{\DIFOincludegraphics}{\includegraphics} %DIF PREAMBLE
\newcommand{\DIFaddincludegraphics}[2][]{{\color{blue}\fbox{\DIFOincludegraphics[#1]{#2}}}} %DIF PREAMBLE
\newcommand{\DIFdelincludegraphics}[2][]{% %DIF PREAMBLE
\sbox{\DIFdelgraphicsbox}{\DIFOincludegraphics[#1]{#2}}% %DIF PREAMBLE
\settoboxwidth{\DIFdelgraphicswidth}{\DIFdelgraphicsbox} %DIF PREAMBLE
\settoboxtotalheight{\DIFdelgraphicsheight}{\DIFdelgraphicsbox} %DIF PREAMBLE
\scalebox{\DIFscaledelfig}{% %DIF PREAMBLE
\parbox[b]{\DIFdelgraphicswidth}{\usebox{\DIFdelgraphicsbox}\\[-\baselineskip] \rule{\DIFdelgraphicswidth}{0em}}\llap{\resizebox{\DIFdelgraphicswidth}{\DIFdelgraphicsheight}{% %DIF PREAMBLE
\setlength{\unitlength}{\DIFdelgraphicswidth}% %DIF PREAMBLE
\begin{picture}(1,1)% %DIF PREAMBLE
\thicklines\linethickness{2pt} %DIF PREAMBLE
{\color[rgb]{1,0,0}\put(0,0){\framebox(1,1){}}}% %DIF PREAMBLE
{\color[rgb]{1,0,0}\put(0,0){\line( 1,1){1}}}% %DIF PREAMBLE
{\color[rgb]{1,0,0}\put(0,1){\line(1,-1){1}}}% %DIF PREAMBLE
\end{picture}% %DIF PREAMBLE
}\hspace*{3pt}}} %DIF PREAMBLE
} %DIF PREAMBLE
\LetLtxMacro{\DIFOaddbegin}{\DIFaddbegin} %DIF PREAMBLE
\LetLtxMacro{\DIFOaddend}{\DIFaddend} %DIF PREAMBLE
\LetLtxMacro{\DIFOdelbegin}{\DIFdelbegin} %DIF PREAMBLE
\LetLtxMacro{\DIFOdelend}{\DIFdelend} %DIF PREAMBLE
\DeclareRobustCommand{\DIFaddbegin}{\DIFOaddbegin \let\includegraphics\DIFaddincludegraphics} %DIF PREAMBLE
\DeclareRobustCommand{\DIFaddend}{\DIFOaddend \let\includegraphics\DIFOincludegraphics} %DIF PREAMBLE
\DeclareRobustCommand{\DIFdelbegin}{\DIFOdelbegin \let\includegraphics\DIFdelincludegraphics} %DIF PREAMBLE
\DeclareRobustCommand{\DIFdelend}{\DIFOaddend \let\includegraphics\DIFOincludegraphics} %DIF PREAMBLE
\LetLtxMacro{\DIFOaddbeginFL}{\DIFaddbeginFL} %DIF PREAMBLE
\LetLtxMacro{\DIFOaddendFL}{\DIFaddendFL} %DIF PREAMBLE
\LetLtxMacro{\DIFOdelbeginFL}{\DIFdelbeginFL} %DIF PREAMBLE
\LetLtxMacro{\DIFOdelendFL}{\DIFdelendFL} %DIF PREAMBLE
\DeclareRobustCommand{\DIFaddbeginFL}{\DIFOaddbeginFL \let\includegraphics\DIFaddincludegraphics} %DIF PREAMBLE
\DeclareRobustCommand{\DIFaddendFL}{\DIFOaddendFL \let\includegraphics\DIFOincludegraphics} %DIF PREAMBLE
\DeclareRobustCommand{\DIFdelbeginFL}{\DIFOdelbeginFL \let\includegraphics\DIFdelincludegraphics} %DIF PREAMBLE
\DeclareRobustCommand{\DIFdelendFL}{\DIFOaddendFL \let\includegraphics\DIFOincludegraphics} %DIF PREAMBLE
%DIF END PREAMBLE EXTENSION ADDED BY LATEXDIFF

\begin{document}

\title{\DIFdelbegin \DIFdel{Koneoppimismallejen }\DIFdelend \DIFaddbegin \DIFadd{Koneoppimismallien }\DIFaddend hyödyntäminen urheilijoiden rasitusvammojen ennustamisessa}
\translatedtitle{in english}
%\studyline{Kaikki opintosuunnat}
\avainsanat{rasitusvammat, koneoppimismallit, ennustaminen}
\keywords{overuse injuries, machine learning, prediction}
\tiivistelma{%

}
\abstract{%

}

\author{Anna Eveliina Kärkkäinen}
\contactinformation{\texttt{anna.e.karkkainen@student.jyu.fi}}
% jos useita tekijöitä, anna useampi \author-komento
\supervisor{Tytti Saksa}
% jos useita ohjaajia, anna useampi \supervisor-komento

% et tarvitse tätä riviä kandidaatin- tai pro gradu -tutkielmassa!
% diplomityössä jätä komento mukaan ja muuta sulkeiden sisään "diplomityö"
\subject{Kandidaatintutkielman}
%\type{tutkimussuunnitelma} 

% vaihda tämä rivi vastaamaan tutkinto-ohjelmaasi, huom. genetiivi ja
% iso alkukirjain
\subject{Tieto- ja ohjelmistotekniikan}

\maketitle

\bigskip


\mainmatter

\chapter{Johdanto}

\DIFaddbegin \DIFadd{Koneoppimisen ala on kasvussa jatkuvasti ja sen mukana myös laskentateho paranee. Laskentatehon kasvaessa koneoppimisen 
hyödyntäminen on lisääntynyt eri tieteenaloilla, mukaan lukien urheilulääketieteessä }\parencite{bhatia_exploring_2024}\DIFadd{. Kasvava huolenaihe on 
nuorten rasitusvammojen yleistyminen, sillä nuorten urheilu on muuttunut yhä kilpailullisemmaksi }\parencite{difiori_overuse_2014}\DIFadd{. 
}

\DIFaddend Tässä kandidaatintutkielmassa \DIFdelbegin \DIFdel{käsitellään urheilijoiden rasitusvammoja ja niiden ennustamista }\DIFdelend \DIFaddbegin \DIFadd{syvennytään aiheeseen kuinka rasitusvammat syntyvät ja kuinka niitä voidaan ennustaa }\DIFaddend koneoppimismallien avulla.
Koneoppimismalleja hyödynnetään jatkuvasti erilaisissa \DIFdelbegin \DIFdel{ennustamisissa. Tässä tutkielmassa paneudutaan siihen}\DIFdelend \DIFaddbegin \DIFadd{ennustamista vaativissa tehtävissä. Tarkoituksena onkin ottaa selvää}\DIFaddend , kuinka urheilijalta 
kerättyä dataa voidaan muuntaa ja hyödyntää koneoppimismalleille sopivaksi. \DIFaddbegin \DIFadd{Tarkempana näkökulmana on juuri nuorten yleistyneet rasitusvammat ja niiden ennustaminen.
}\DIFaddend 

\DIFdelbegin \DIFdel{Tutkimuksen aiheena on ottaa selvää, kuinka urheilijoiden rasitusvammat syntyvät sekä kuinka niitä voidaan ennustaa koneoppimismallejen 
avulla.
Tarkempi näkökulma on selvittää, kuinka tehokkaita koneoppimismallit ovat ennustamaan rasitusvammoja. }\DIFdelend Aiheesta suoraan ei ole kovinkaan paljoa tutkimusta, joten \DIFdelbegin \DIFdel{koen tutkimuskysymyksen olevan }\DIFdelend \DIFaddbegin \DIFadd{tutkimuskysys on }\DIFaddend ajankohtainen ja mielenkiintoinen. \DIFdelbegin %DIFDELCMD < 

%DIFDELCMD < %%%
\DIFdel{Juuri kyseinen aihe valikoitui, sillä olen kiinnostunut suuntaamaan opintoni sekä tulevaisuuden työni juuri urheilun ja 
tekniikan yhdistävälle alalle. Koen myös, ettei aiheesta ole kovinkaan paljoa selkeää ja yksinkertaistettua tietoa saatavilla 
suomeksi. Tarkoituksen on }\DIFdelend \DIFaddbegin \DIFadd{Tutkielman tarkoituksen on siis }\DIFaddend koota erilaisia artikkeleita ja 
tutkimuksia yhteen, jotta \DIFdelbegin \DIFdel{saan }\DIFdelend \DIFaddbegin \DIFadd{lukija saa }\DIFaddend monipuolisen ja yksinkertaistetun kokonaisuuden rasitusvammojen syntymisestä sekä koneoppimismalleilla ennustamisesta.


\chapter{\DIFdelbegin \DIFdel{Tutkimusmenetelmät}\DIFdelend \DIFaddbegin \DIFadd{Käsitteet}\DIFaddend }

\DIFdelbegin \DIFdel{Valitsemani aiheen osa-alueista löytyy kattavasti eri muotoista aineistoa, jota voin hyödyntää kirjallisuuskartoituksen kirjoittamisessa. 
Aiheesta suoraan ei ole löytynyt kuin muutamia tutkimuksia. Olen löytänyt esimerkiksi hyvin lajispesifisiä tutkimuksia, mutta suoraan yleis-
kattavaa ja kartoittavaa tutkimusta en ole löytänyt. Muun muassa näiden syiden takia päädyin valitsemaan tämän aiheen kandidaatintutkielmaani 
varten.
}%DIFDELCMD < 

%DIFDELCMD < %%%
\DIFdel{Tällä hetkellä lähteitä olen hakenut ja löytänyt lähinnä Keenioukselta, Scopucsesta ja IEEE Xploresta. Etsiessäni lähteitä koitan valikoida lähteet sen 
perusteella ovatko ne vertaisarvioituja, ovatko ne julkaistu luotettavalla palstalla sekä onko lähteisiin viitattu useaan otteeseen. 
Hakusanoina olen käyttänyt muun muassa "machine learning AND overuse injuries", "different types of machine learning"sekä "injury prevention AND machine learning".
}%DIFDELCMD < 

%DIFDELCMD < %%%
\chapter{\DIFdel{Käsitteet}}
%DIFAUXCMD
\addtocounter{chapter}{-1}%DIFAUXCMD
%DIFDELCMD < 

%DIFDELCMD < %%%
\DIFdelend \section{Rasitusvammat}

\DIFdelbegin \DIFdel{(tähän tarkoitus etsiä lähde, joka kertoo mitä rasitusvammat ovat)
}%DIFDELCMD < 

%DIFDELCMD < %%%
\DIFdelend \DIFaddbegin \DIFadd{Rasitusvammat syntyvät ajan myötä, toisin kuin akuutit vammat, jotka syntyvät suoraan tapahtumasta. }\DIFaddend Rasitusvammoja voi syntyä, kun yksilö altistuu 
toistuvalle rasitukselle. Kun kehon kudokset kokee liiallista rasitusta, liian vähäistä palautumista tai molempia, kehon kyky palautua ylittyy \parencite{johnson_overuse_2008}.
Tarkemmin rasitusvammat ovat seurausta luuhun, lihakseen ja/tai jänteeseen kohdistuvasta kumulatiivisesta 
mikrotraumasta \parencite{brenner_overuse_2024}.

\textcite{johnson_overuse_2008} toteaa rasitusvammojen syyt jakautuvan yleensä kahteen \DIFdelbegin \DIFdel{osaa}\DIFdelend \DIFaddbegin \DIFadd{osaan}\DIFaddend . Sisäisiin tekijöihin \DIFdelbegin \DIFdel{(}\DIFdelend \DIFaddbegin \DIFadd{eli }\DIFaddend yksilöstä johtuviin \DIFdelbegin \DIFdel{) }\DIFdelend ja ulkoisiin 
tekijöihin \DIFdelbegin \DIFdel{(}\DIFdelend \DIFaddbegin \DIFadd{eli }\DIFaddend ympäristöstä johtuviin\DIFdelbegin \DIFdel{)}\DIFdelend . Hän mainitsee sisäisiä tekijöitä olevan muun muassa anatomiset virheasennot, heikko fyysinen kunto 
sekä aikaisemmat vammat. \DIFdelbegin \DIFdel{Puolestaan isäisiin tekijöihin ei henkilöllä ole yleensä vaikutusvaltaa. Ulkoisiin tekijöihin yksilöllä 
on kuitenkin enemmän vaikutusvaltaa ja näitä ovat , }\DIFdelend \DIFaddbegin \DIFadd{Ulkoisia tekijöitä puolestaan ovat }\DIFaddend esimerkiksi väärät harjoittelumetodit, huonot tekniikat sekä 
valmentajien tai harjoittelu kavereiden aiheuttama paine. \DIFdelbegin %DIFDELCMD < \parencite{johnson_overuse_2008}
%DIFDELCMD < %%%
\DIFdelend \DIFaddbegin \DIFadd{Ulkoisiin tekijöihin yksilön on mahdollista vaikuttaa, toisin kuin sisäisiin tekijöihin.
}\DIFaddend 

\DIFaddbegin \DIFadd{Nuorten kilpailullinen urheilu on muuttunut yhä vaativammaksi, mikä on johtanut rasitusvammojen yleistymiseen }\parencite{weerasinghe_computer_2016}\DIFadd{.
}\DIFaddend \section{Koneoppimismallit}

\DIFdelbegin \subsection{\DIFdel{Mitä on koneoppimismallit}}
%DIFAUXCMD
\addtocounter{subsection}{-1}%DIFAUXCMD
%DIFDELCMD < 

%DIFDELCMD < %%%
\DIFdel{Koneoppiminen (eng}\DIFdelend \DIFaddbegin \DIFadd{Koneoppiminen (engl}\DIFaddend . Machine Learning) tarkoittaa automatisoitua prosessia, jossa datasta tunnistetaan säännönmukaisuuksia 
tehtävien, kuten luokittelun ja ennustamisen, suorittamiseksi \parencite{black_introduction_2023}. Koneoppimismallit jaetaan 
yleisesti kolmeen eri kategoriaan: ohjattu oppiminen, ohjaamaton oppiminen ja vahvistettu oppiminen. 

%DIF < \Näistä ohjattua oppimista käytetään yleisesti ennustamiseen, jonka vuoksi keskityn siihen tässä tutkielmassa. \parencite{singh_machine_2021}
\DIFaddbegin \DIFadd{Tässä työssä käsiteltävät menetelmät rajautuvat ohjattuun oppimiseen, sillä ohjattun oppimisen menetelmät ovat käytetyimpiä ennsutamiseen }\parencite{singh_machine_2021}\DIFadd{.
}\DIFaddend 

\subsection{Ohjattu oppiminen}

\DIFdelbegin \DIFdel{Ohjatun oppimisen luokittelualgoritmeja ovat muun muassa lähimmän }\DIFdelend \DIFaddbegin \DIFadd{Ohjatussa oppimisessa mallille syötetään dataa, jonka lopputulokset tunnetaan. Malli oppii tunnistamaan datasta säännönmukaisuuksia ja oppii näin luokittelemaan tai ennustamaan 
havaitsemattomia lopputuloksia uudesta datasta }\parencite{black_introduction_2023}\DIFadd{. Ohjatun oppimisen menetelmiä on useita, joista tässä työssä merkittävimpiä ovat seuraavat:
}


\begin{tabular}{|p{7cm}|p{8cm}|}
\hline
\DIFadd{Algoritmi }& \DIFadd{Selitys }\\ \hline
\DIFadd{Lähimmän }\DIFaddend naapurin menetelmä (\DIFdelbegin \DIFdel{eng}\DIFdelend \DIFaddbegin \DIFadd{engl}\DIFaddend . K-Nearest Neighbors, K-NN) \DIFdelbegin \DIFdel{, päätöspuu, satunnaismetsä }\DIFdelend \DIFaddbegin & \DIFadd{Algoritmia käytetään }\DIFaddend sekä \DIFdelbegin \DIFdel{tukivektorikoneet (eng}\DIFdelend \DIFaddbegin \DIFadd{regressio- että luokitteluongelmien ratkaisemiseen. 
Käytetään ennustamiseen samankaltaisten tapausten keskiarvon ja mediaanin perusteella sekä laskemaan tulosdata perustuen K-määrän samankaltaisimpien 
tapausten esiintymistiheyteen }\parencite{singh_machine_2021}\DIFaddend . \DIFaddbegin \\ \hline
\DIFadd{Päätöspuu (engl. Decision trees) }& \DIFadd{Data jaetaan yhä samankaltaisempiin ryhmiin perustuen muuttujaan, joka maksimoi tuloksena olevien ryhmien samankaltaisuuden. Uudet luokittelut ja ennusteet tehdään arvioimalla uutta yksilöä näiden jakokriteerien mukaisesti }\parencite{black_introduction_2023}\DIFadd{.}\\ \hline
\DIFadd{Satunnaismetsä (engl. Random forest) }& \DIFadd{Päätöspuihin perustuva ryhmämenetelmä, joka yhdistää useita itsenäisesti satunnaistettuja päätöspuita. Jokainen puu rakennetaan hyödyntämällä satunnaista otantaa ja satunnaisia muuttujajoukoja, mikä varmistaa mallin vakauden ja vähentää virhealttiutta }\parencite{breiman_random_2001}\DIFadd{. }\\ \hline
\DIFadd{Tukivektorikoneet (engl. }\DIFaddend Support vector machines, SVM) \DIFdelbegin \DIFdel{. Luokittelu ja ennustaminen ovat ohjatun oppimisen 
yleisimpiä käyttötarkoituksia }\DIFdelend \DIFaddbegin & \DIFadd{Malli jakaa havaintoaineiston osiin n-ulotteisen hypertason perusteella. Uudet yksilöt luokitellaan sen mukaan, minkä rajatun alueen sisälle heidän ominaisuutensa osuvat }\DIFaddend \parencite{black_introduction_2023}.\DIFdelbegin %DIFDELCMD < 

%DIFDELCMD < %%%
\chapter{\DIFdel{Rasitusvammojen ennustaminen koneoppimismalleilla}}
%DIFAUXCMD
\addtocounter{chapter}{-1}%DIFAUXCMD
%DIFDELCMD < 

%DIFDELCMD < %%%
\section{\DIFdel{Käytetyimmät koneoppimismallit ennustamiseen}}
%DIFAUXCMD
\addtocounter{section}{-1}%DIFAUXCMD
%DIFDELCMD < 

%DIFDELCMD < %%%
\DIFdel{\textcite{black_introduction_2023} kuvaavat artikkelissaan taulukkona yleisimpiä koneoppimisen lähestymistapoja luokitteluun ja ennustamiseen.
Taulukosta löytyy jo aijemmin mainitsemani päätöspuu, satunnaismetsä, tukivektorikoneet sekä lähimmän naapurin menetelmä}\DIFdelend \DIFaddbegin \\ \hline
\DIFadd{Neuroverkot (engl}\DIFaddend . \DIFdelbegin \DIFdel{Näiden lisäksi 
taulukossa mainitaan muun muassa neuroverkot (eng. }\DIFdelend Neural networks) \DIFdelbegin \DIFdel{, lineaarinen }\DIFdelend \DIFaddbegin & \DIFadd{- }\\ \hline
\DIFadd{Lineaarinen }\DIFaddend ja logistinen regressio \DIFdelbegin \DIFdel{, }\DIFdelend \DIFaddbegin & \DIFadd{- }\\ \hline
\DIFaddend Naiivi Bayes (\DIFdelbegin \DIFdel{eng}\DIFdelend \DIFaddbegin \DIFadd{engl}\DIFaddend . Naive Bayes) \DIFdelbegin \DIFdel{sekä 
malliyhdistelmät (eng}\DIFdelend \DIFaddbegin & \DIFadd{- }\\ \hline
\DIFadd{Malliyhdistelmät (engl}\DIFaddend . Ensembles) \DIFdelbegin \DIFdel{. 
}\DIFdelend \DIFaddbegin & \DIFadd{- }\\ \hline
\end{tabular}
\DIFaddend 


\DIFdelbegin \DIFdel{Samat koneoppimismallit mainitaan myös \textcite{musat_diagnostic_2024} artikkelissa. Kyseinen artikkeli keskittyy tekoälyn hyödyntämiseen 
loukkaantumisen ennustamisessa, mutta käy tiiviisti läpi myön koneoppimismalleja, joita voidaan hyödyntää ennustamisessa. }\DIFdelend \DIFaddbegin \chapter{\DIFadd{Rasitusvammojen ennustaminen koneoppimismalleilla}}
\DIFaddend 

\DIFaddbegin \section{\DIFadd{Käytetyimmät koneoppimismallit ennustamiseen}}

\DIFadd{- Kysymys: miten tämän kappaleen saa järkevästi kytkettyä edelliseen kappaleeseen "ohjattu oppiminen"? Onko yllä oleva taulukko hyvä idea?
}


\DIFaddend \section{Urheiluvammojen ennustaminen}

\subsection{Datan keruu ja sen esikäsittely}

\DIFaddbegin \DIFadd{Datan esikäsittely on tärkeää, sillä koneoppimismallit vaativat datan olevan tietyssä muodossa. Monesti urheilijalta kerätty data on esimerkiksi saatua 
dataa urheilukellosta tai kyselyiden pohjalta saatua dataa. Tästä johtuen dataa on esikäsiteltävä, jotta malli osaa sitä hyödyntää. Saatua dataa muun muassa puhdistetaan, 
normalisoidaan ja muunnetaan eri muotoon }\parencite{murugan_comparison_2025}\DIFadd{. Esikäsittelyyn kuuluu myös esimerkiksi datan imputointi eli paikkaus puuttuvien arvojen 
hallitsemiseksi }\parencite{yuan_machine_2026}\DIFadd{. 
}

\DIFadd{Rasitusvammoja tutkiessa yleisin datan keruu muoto on erilaiset kyselyt, joihin urheilijat vastaavat. }\DIFaddend Tutkimuksessa \parencite{murugan_comparison_2025} 
vertailtiin koneoppimismalleja urheilijoiden loukkaantumisten ennustamiseen. Tutkimuksessa dataa kerättiin useilta urheilijoilta määritellyn ajanjakson aikana. 
Jokainen aineiston merkintä vastasi yhtä harjoitus päivää ja sisälsi muun muassa seuraavat ominaisuudet: nr.sessions (urheilijan yhden päivän aikana suorittamien 
harjoitusten kokonaismäärä), total km (kaikkien harjoitusten aikana kertynyt kokonaiskilometrimäärä), km z3-4 (sykealueilla 3 ja 4 taitettu matka) sekä strength 
training (päivän aikana tehty voimaharjoittelun kokonaismäärä). 

\textcite{de_leeuw_personalized_2022} tutkimuksessa puolestaan urheilijat vastasivat päivittäin OSTRC-kyselyyn (\DIFdelbegin \DIFdel{eng}\DIFdelend \DIFaddbegin \DIFadd{engl}\DIFaddend . Oslo Sports Trauma 
Research Center questionnaire). 

\subsection{Algoritmit}

- algoritmit, joita lähteistä löydettiin: Subgroup discovery \parencite{de_leeuw_personalized_2022}. Satunnaismetsä, logistinen regressio ja kNN \parencite{murugan_comparison_2025}

\section{Lopputulokset}

- miten paljon koneoppimismalleja hyödynnetty rasitusvammojen ennustamisessa

- onko niistä mitä hyötyä

\chapter{Yhteenveto}

Koneoppimismalleja käytetään jatkuvasti erilaisissa ennustamisissa. Tässä tutkielmassa paneuduttiin siihen, kuinka urheilijalta kerättyä 
dataa voidaan muuntaa ja hyödyntää koneoppimismalleille sopivaksi.

Lopussa yhteenvedossa tulee ilmi tiivis kokonaisuus juuri kertomistani asioista. Yhteenvedossa tulee näkymään selkeä ymmärrykseni aiheesta.

\chapter{Johtopäätökset}

Tutkielaman aiheesta ei kovin paljoa tutkimusta löydy. Aihe on kuitenkin hyvin tuore ja varmasti lisää tutkimusta tulevaisuudessa tehdään. 
Koneopimismallit kehittyvät jatkuvasti ja sitä myöten tutkimukset siitä kuinka niitä voidaan ennustamiseen hyödyntää, lisääntyvät.

\printbibliography


\end{document}
